\section{\texorpdfstring{$\mathsf{TS}\left(M\right)$}{TS(M)}和\texorpdfstring{$\mathsf{S}\left(M\right)$}{S(M)}的关系}

\begin{proposition}{}{}
	设$s:\mathsf{T}\left(M\right)\hookrightarrow\mathsf{TS}\left(M\right)$由$M\hookrightarrow\mathsf{TS}\left(M\right)$诱导,则存在唯一的分次$A$-代数同态$\varphi_{M}:\mathsf{S}\left(M\right)\to\mathsf{TS}\left(M\right)$使下图交换.
	\begin{center}
		\begin{tikzcd}
			& {\mathsf{T}\left(M\right)} & \\
			{\mathsf{S}\left(M\right)} && {\mathsf{TS}\left(M\right)}
			\arrow["\rho"', two heads, from=1-2, to=2-1]
			\arrow["s", hook, from=1-2, to=2-3]
			\arrow["{\exists!\varphi_{M}}"', dashed, from=2-1, to=2-3]
		\end{tikzcd}
	\end{center}
	此外,记$i:\mathsf{TS}\left(M\right)\hookrightarrow\mathsf{T}\left(M\right)$为典范映射,定义分次$A$-线性映射$\psi_{M}\coloneq\rho\circ i$.
\end{proposition}

\begin{proposition}{}{}
	设$u\in\Hom_{A}\left(M,N\right)$,则如下图表交换.
	\begin{center}
		\begin{tikzcd}
			{\mathsf{S}\left(M\right)} & {\mathsf{TS}\left(M\right)} & {\mathsf{S}\left(M\right)} \\
			{\mathsf{S}\left(N\right)} & {\mathsf{TS}\left(N\right)} & {\mathsf{S}\left(N\right)}
			\arrow["{\varphi_{M}}", from=1-1, to=1-2]
			\arrow["{\mathsf{S}\left(u\right)}"', from=1-1, to=2-1]
			\arrow["{\varphi_{N}}", from=1-2, to=1-3]
			\arrow["{\mathsf{TS}\left(u\right)}"', from=1-2, to=2-2]
			\arrow["{\mathsf{S}\left(u\right)}", from=1-3, to=2-3]
			\arrow["{\varphi_{N}}"', from=2-1, to=2-2]
			\arrow["{\psi_{N}}"', from=2-2, to=2-3]
		\end{tikzcd}
	\end{center}
\end{proposition}

\begin{proposition}{}{}
	设$\left(M_{\lambda}\right)_{\lambda\in\Lambda}$是$M$的子模族,$M=\bigoplus\limits_{\lambda\in\Lambda}M_{\lambda}$,则如下图表交换(竖直方向的箭头为典范映射).
	\begin{center}
		\begin{tikzcd}
			{\bigotimes\limits_{\lambda\in\Lambda}\mathsf{S}\left(M_{\lambda}\right)} && {\bigotimes\limits_{\lambda\in\Lambda}\mathsf{TS}\left(M_{\lambda}\right)} \\
			{\mathsf{S}\left(M\right)} && {\mathsf{TS}\left(M\right)}
			\arrow["{\bigotimes\limits_{\lambda\in\Lambda}\varphi_{M_{\lambda}}}", from=1-1, to=1-3]
			\arrow[from=1-1, to=2-1]
			\arrow[from=1-3, to=2-3]
			\arrow["{\varphi_{M}}"', from=2-1, to=2-3]
		\end{tikzcd}
	\end{center}
\end{proposition}

\begin{proposition}{自由模上的双代数同态}{}
	若$M$是自由模,则$\varphi_{M}$是分次双代数同态.
\end{proposition}

\begin{proposition}{}{}
	固定$n\in\N$.
	\begin{enumerate}
		\item 若$u\in\mathsf{S}^{n}\left(M\right)$,则$\psi_{M}\left(\varphi_{M}\left(u\right)\right)=n!u$.
		\item 若$v\in\mathsf{TS}^{n}\left(M\right)$,则$\varphi_{M}\left(\psi\left(v\right)\right)=n!v$.
	\end{enumerate}
\end{proposition}

\begin{corollary}{特征$0$域上对称代数与对称张量代数同构}{}
	设$A$是$\Q$-代数,则典范映射$\mathsf{S}\left(M\right)\to\mathsf{TS}\left(M\right)$是代数同构.进一步,若$M$是自由模,该同构为分次双代数同构.
\end{corollary}

\begin{corollary}{$\Q$-代数上$n$次对称张量代数的生成元}{}
	设$A$是$\Q$-代数,则对每个$n\in\N$,$\mathsf{TS}^{n}\left(M\right)$由$M$在$\mathsf{TS}\left(M\right)$中的元素的$n$次幂生成.
\end{corollary}








































